% ============================================================
% KI-Verzeichnis – automatisierte Einbindung via \nocite{*}
% ============================================================
%
% Statt die Einträge manuell aufzulisten:
%
%   \nocite{prompt01, prompt02, ..., promptArchitectureGen}   ← VERALTET / fehleranfällig
%
% wird hier \nocite{*} verwendet. Damit werden ALLE Einträge
% der .bib-Datei geladen. Der keyword-Filter in
% \printbibliography stellt sicher, dass nur Einträge mit
% keywords={ki} im Verzeichnis erscheinen.
%
% Neue Prompts müssen nur noch in ki_prompts.bib eingetragen
% werden – kein manuelles Nachtragen in \nocite nötig.
% ============================================================

% Alle Einträge aus ki_prompts.bib laden (automatisch):
\nocite{*}

% KI-Verzeichnis ausgeben (gefiltert nach keyword=ki):
\printbibliography[
  title={KI-gestützte Werkzeuge (Prompts)},
  heading=bibnumbered,
  keyword=ki,
  resetnumbers=true
]
